\uuid{rhzT}
\titre{Quotient/DL usuels/DL en l'infini}
\auteur{hanani ; mzari}
\organisation{exo7}
\licence{CC BY-NC-SA 4.0}
\datecreate{2020-01-20}
\niveau{L1}
\module{Analyse}
\chapitre{Développement limité}
\sousChapitre{Calculs}
\difficulte{3}
\theme{Opérations sur les DL, Quotient, DL usuels, DL en l'infini}
\exo7qcmid{lille-2.developpements-limites.029}

\contenu{
\texte{
Soit $a\in \Rr$ et $f$ une fonction. On dit que $f(x)$ est un $o\left((x-a)^n\right)$ au voisinage de $a$ et on écrit $f(x)=o((x-a)^n)$ si $\displaystyle \lim _{x\to a}\frac{f(x)}{(x-a)^n}=0$. On note par $DL_n(a)f(x)$ le développement limité de $f$ à l'ordre $n$ en $a$. On note aussi $DL_n(a^+)f(x)$ (resp. $DL_n(a^-)f(x)$) le développement limité de $f$ à l'ordre $n$ à droite en $a$ (resp. à gauche en $a$).
}

\begin{question}
\qclassification{125.02}
\qtags{motcle=Quotient/DL usuels/DL en l'infini}
Soient $f$ et $g$ telles que $\displaystyle f(x)=\frac{x^2}{x-1}\ln\left(\frac{x+1}{x}\right)$ et $\displaystyle g(t)=f\left(\frac{1}{t}\right)$. Parmi les affirmations suivantes, cocher celles qui sont vraies :
\begin{answers}
\good{$\displaystyle g(t)=\frac{\ln (1+t)}{t-t^2}$.}
\bad{$\displaystyle DL_2(0)g(t)=1+\frac{t}{2}+\frac{t^2}{2}+o(t^2)$.}
\bad{$\displaystyle DL_2(0)(+\infty)f(x)=1+\frac{1}{x}+\frac{1}{2x^2}+o\left(\frac{1}{x^2}\right)$.}
\good{$\displaystyle DL_2(0)(+\infty)f(x)=1+\frac{1}{x}+\frac{5}{6x^2}+o\left(\frac{1}{x^2}\right)$.}
\end{answers}
\begin{explanations}
D'abord $\displaystyle g(t)=\frac{\ln (1+t)}{t-t^2}=1+\frac{t}{2}+\frac{5t^2}{6}+o(t^2)$. Donc
$$DL_2(+\infty)f(x)=1+\frac{1}{x}+\frac{5}{6x^2}+o\left(\frac{1}{x^2}\right).$$
\end{explanations}
\end{question}
}
